\documentclass[10pt,a4paper]{amsart}
\usepackage[margin=19mm]{geometry}
\usepackage{amsmath,amssymb,mathtools}
\usepackage[scr=rsfs,cal=euler]{mathalfa}
\usepackage{xfrac}
\usepackage[unicode,hidelinks,bookmarks=false]{hyperref}
\newcommand{\ie}{i.e.\ }
\newcommand{\cf}{cf.\ }
\newcommand{\resp}{respectively\ }
\newcommand{\Subsection}{\subsection}
\providecommand{\nobreakdash}{\nobreak\hbox{-}}
\setlength{\emergencystretch}{2em}
\multlinegap0pt
\usepackage{bm}
\usepackage{bbm}
\usepackage{dsfont}
\usepackage{xspace}
\usepackage{cancel}
\usepackage{mathtools}
\usepackage{amsthm}
\makeatletter
\@ifundefined{thm}{\newtheorem{thm}{Thm}}{}
\@ifundefined{lem}{\newtheorem{lem}[thm]{Lemma}}{}
\makeatother
\def\wt{\mathrm{wt}}  % weight 
\let\vep\varepsilon

\def\sea{{\scriptscriptstyle\searrow}}
\title[On some components of $L(\rho)\otimes L(\rho)$ associated wi]{On some components of $L(\rho)\otimes L(\rho)$ associated with rooted trees for symmetrizable Kac-Moody algebras}
\author{Rekha Biswal, Patrick Polo}
\date{Source: arXiv:2606.27197v2 (CC BY 4.0)}
\begin{document}
\maketitle

\noindent\emph{Source and licence.} On some components of $L(\rho)\otimes L(\rho)$ associated with rooted trees for symmetrizable Kac-Moody algebras. Version arXiv:2606.27197v2 (CC BY 4.0), distributed under the Creative Commons Attribution 4.0 International licence, as recorded in the arXiv licence metadata for this exact version and shown on the abstract page https://arxiv.org/abs/2606.27197v2. Full rights notice: licenses/arxiv-2606.27197-CC-BY-4.0.txt.\par\smallskip

\noindent\textit{Editorial scope.} Counted unit: the Lem whose source label is \texttt{lem-pi-I} in the article (source0.tex lines 346--349, source label \texttt{lem-pi-I}), delivered together with the article's own complete original proof of it (source0.tex lines 351--379, one \texttt{proof} environment of 29 lines). No mathematics is written, repaired or re-derived: statement and proof are the article's own text line for line. Required context retained uncounted: lem-BS (lines 246--267). Every definition, display and label of those lines is retained unchanged. Omitted from this package and not redistributed: 1--245, 268--345, 350--350, 380--1293. Nothing inside a retained excerpt is omitted or altered: every retained body is delivered exactly as the article's lines are written, so no omission receipt of any kind accompanies this package. Citations: the retained excerpts carry no citation command at all, so no bibliography entry of the article is transcribed and no reference is redistributed. Reference adaptation: none was needed and none was made; every reference inside the delivered unit resolves inside it, so no unresolved reference or citation is produced. Printed slips kept literal: no printed word, symbol, hypothesis or number of the counted statement or of its proof was altered. Layout and class: the article's own class is \texttt{amsart} with packages including babel, xy, amssymb, tikz; the wrapper is the lane's pinned \texttt{amsart} editorial wrapper at 10pt on A4, which restates the theorem-like environments the retained text needs and adds the article's own macro definitions that the retained text needs (\texttt{\textbackslash vep}, \texttt{\textbackslash wt}), with extra wrapper packages (bm, bbm, dsfont, xspace, cancel, mathtools). The delivered numbering is the unit's own. Assets: no retained span contains a figure environment, a graphics-inclusion command, a PDF-page inclusion, a table or a TikZ picture; no image file is copied into this package, the package declares no asset dependency and no image file is redistributed with it. Licence: version arXiv:2606.27197v2 of this article is distributed under the Creative Commons Attribution 4.0 International licence, as recorded in the arXiv licence metadata for this exact version and shown on the abstract page https://arxiv.org/abs/2606.27197v2. A full rights notice is delivered at licenses/arxiv-2606.27197-CC-BY-4.0.txt. Characters: every retained excerpt is pure ASCII, so the delivered engine, XeLaTeX with its default Latin Modern text font, reports no missing character.
\begin{lem}\label{lem-BS} 
Let $x_j \in B_j$, for crystals $B_1, \dots, B_n$. Consider the element $x = x_1 \otimes \cdots \otimes x_n$ of 
$B_1 \otimes \cdots \otimes B_n$. Then, for all $s\in \Delta$, 

\smallskip {\rm (a)} $\varphi_s(x)$ is given by 
$\varphi_s(x) = \max_{1\leq j\leq n} \Big( \varphi_s(x_j) + \Sigma_s(x,j) \Big)$ 
and if  $j$ is the first index where the maximum is attained, then
$$
f_s(x)
=
x_1 \otimes \cdots \otimes f_s(x_j)\otimes \cdots \otimes x_n.
$$

{\rm (b)} $\vep_s(x)$ is given by 
$\vep_s(x) = \max_{1\leq j\leq n} \Big( \vep_s(x_j) + \Sigma'_s(x,j) \Big)$
and if  $j$  is the last index where the maximum is attained, then
$$
e_s(x)
=
x_1 \otimes \cdots \otimes e_s(x_j)\otimes \cdots \otimes x_n.
$$
\end{lem} 

\begin{lem}\label{lem-pi-I} 
One has $\pi_I = \Big( \bigotimes_{i=1}^N f_i \pi_i \Big) \otimes \Big( \bigotimes_{j > N} \pi_j \Big)$, recalling that 
$N = \vert I \vert$. 
\end{lem}

\begin{proof} By induction on $\vert I\vert$. There is nothing to prove if $I = \varnothing$, so we may assume 
$I \not= \varnothing$
and the result proved for $I' = I - \{N\}$, where $N$ is the largest element of $I$. 
Set $X' = \pi_{I'}$. We want to compute $f_N X'$. 

\smallskip 
By the inductive hypothesis, one has $x_j = \pi_j$ if $j\geq N$ and $x_j = f_j \pi_j$ if $j < N$. 
From this it follows that $\varphi_N(x_j) = \delta_{N,j} = \wt_N(x_j) $ for $j\geq N$, while for $j < N$ one has: 
$\varphi_N(x_j) = 1 = \wt_N(x_j)$ if $j = N^-$ and $\varphi_N(x_j) = 0 = \wt_N(x_j)$ else. 

Thus, if $N$ has a predecessor $N^-$ one has: 
$$
\Phi_N(x,j) = 
\begin{cases} 
0 & \text{if } j < N^-, 
\\
1 & \text{if } N^- \leq j < N, 
\\
2 & \text{if } j \geq  N,
\end{cases}
$$
and otherwise one has $\Phi_N(x,j) = 0$ if $j < N$ and $\Phi_N(x,j) = 1$ if $j\geq N$. 
In both cases, $N$ is the first index where the maximum of $\Phi_N(x,j)$ is attained. Hence, by Lemma \ref{lem-BS} (a), 
one obtains, with obvious notation: 
$$
f_N X' = X'_{\leq N-1} \otimes f_N \pi_N \otimes X'_{> N}.
$$
This proves the lemma.
\end{proof} 

\end{document}
